\begin{itemize}
    \item \textbf{Cross-section: no interpretable estimate.} Across seeds the point estimates move enormously; for example, a random-forest \(\hat p\) with a kernel \(f\) gives between \(0.02\) and \(26.3\), and a net \(\hat p\) with a kernel \(f\) between \(0.6\) and \(13.2\) (\Cref{tab:seed9}). The main-run values in \Cref{tab:cross9} are therefore one draw from a wide distribution, and the bootstrap intervals, which hold the tuning parameters and the random folds fixed, do not reflect that variation. With \(n = 101\) and a first stage whose out-of-fold \(R^2\) is between \(0.01\) and \(0.10\) (gain from the instrument \(0.03\) to \(0.06\)), the data do not determine \(\tilde f\), and the estimates mostly reflect the cross-validation folds and the resampled draws.
    \item \textbf{Panel: more stable, and the class of \(f\) matters more than the class of \(\hat p\).} The three first stages are similar (out-of-fold \(R^2\) \(0.14\) to \(0.16\), each gain significant at the \(0.02\) level). The random-forest \(f\) gives small estimates in every cell, \(0.7\) to \(2.1\) across all seeds and \(\hat p\) classes, consistent with the shrinkage seen earlier. The kernel \(f\) gives \(7.3\) to \(10.1\) across seeds when paired with the random-forest or net \(\hat p\), against a linear 2SLS estimate of \(12.34\) (s.e.\ \(4.35\)); the bootstrap intervals for those cells contain it. The kernel \(f\) with the kernel \(\hat p\) is less stable (\(1.9\) to \(9.5\) across seeds), and the net \(f\) varies from \(1.5\) to \(9.9\) across seeds and \(\hat p\) classes.
    \item The seed reruns used 5 (cross-section) and 3 (panel) additional seeds and no bootstrap, so the ranges in \Cref{tab:seed9} understate the true spread.
    \item In the panel, the log-density gain from the instrument is significant only for the kernel \(\hat p\) at the \(10\%\) level (\(p = 0.098\)), not for the forest or net (\(p = 0.51\), \(0.33\)): the instrument improves the prediction of the mean of \(A\) but not clearly the shape of the residual distribution, and we have not explained why the two measures disagree.
    \item All of this rests on our location-model assumption for \(\hat p\) (an additive error independent of \((Z, X)\)), the choice of \(M = 50\) draws, and tuning on the same folds used to report fit. We do not draw a conclusion about the causal effect of immigration from these estimates.
\end{itemize}
